\documentclass[10pt,english]{article}
\usepackage{babel,amsmath,amsfonts,amssymb,color,fancybox}
\usepackage[utf8]{inputenc}
\usepackage[a4paper, inner=2.5cm, outer=2.5cm, top=2.0cm, bottom=2cm, bindingoffset=0cm]{geometry}
\usepackage{graphicx}
\usepackage{stmaryrd}
\usepackage{amsmath}
\usepackage{hyperref}
\usepackage{float}
%\usepackage{times}                  % stile delle sections

\newcommand{\reals}{\mathbf{R}}
\parindent=0pt \parskip=3pt
%
\title{Review of Research Paper number 24}
\date{version: \today}%Meeting held on: April the 2nd}
\author{Lucien Le Gall, Noé Casteleyn, Loris Manganelli}



\begin{document}

\maketitle
%\small

\section{Introduction}

The Python code for Section~\ref{sec:practice} is available in the file \lstinline{ADMM\_FHE.py}, and its execution is provided in the accompanying notebook. The project repository can be accessed at:

\begin{center}
\url{https://github.com/Loris-Manganelli/projet-cooperative-optimisation/tree/main}
\end{center}

\subsection{Summary of the paper}

% Describe in your own words [10 lines max]

The article \cite{Binfet_2023} focuses on privacy preserving \textit{via} encryption (which is one of three ways
among anonymisation, encryption, and differential privacy to ensure data privacy). It mainly says that existing approaches of encryption
are cloud-based solution. Thus, the authors have the original idea to combine peer-to-peer (what they call cooperative)
ADMM with encryption (based on a fully homomorphic encryption (FHE) scheme).

\subsection{Problem Formulation}
The article focuses on the following problem:
\begin{equation*}
\begin{aligned}
& \min_{z_i} \quad 
\sum_{i=1}^M f_i(z_i,p_i)
\text{ s.t. }& z_i=(\zeta)_{\kappa_i}
\end{aligned}
\end{equation*}

where $M$ is the number of agents, and $\zeta$ a variable that enables consensus. The function $f_i$ is defined as:
\begin{equation}
f_i(z_i, p_i) =
\frac{1}{2} z_i^T H_i z_i + p_i^T F_i^T z_i +
\begin{cases}
0, & \text{if } G_i z_i = E_i p_i,\\
\infty, & \text{otherwise.}
\end{cases}
\end{equation}

The affine constraint ensure that the dynamic is feasible, it is a link between the variables and the states (useful in robotic applications).

The articles introduces the following notation: $z_i=\begin{pmatrix}
\alpha_i\\
\gamma_i
\end{pmatrix}$, and $\rho_i=\begin{pmatrix}
\beta_i\\
\delta_i
\end{pmatrix}$. The variable $\alpha_i$ is the agent's variable, while $\gamma_i$ represents other information the agent have over its neighbors. 

We can adapt the course main problem
\begin{equation}\label{problème du cours}
    \min_{x\in \mathbb R^n} \sum_{i=1}^N f_i(x)
\end{equation}

with the following formulation: 
\begin{equation}\label{main problem}
    \begin{aligned}
    &\min_{\alpha_i\in\mathbb{R}^{m}, \gamma_{ij}\in\mathbb{R}^{m}} \sum_{i=1}^{N} f_i(\alpha_i)\\
    \text{s.t } & \gamma_{ij}=\zeta_{j} \quad \forall i\sim j & \textit{(i)}\\ 
    & \alpha_i=\zeta_i \quad \forall i \in M & \textit{(ii)}\\
    & \alpha_i=\gamma_{ij} \quad \forall i\sim j & \textit{(iii)}
    \end{aligned}
\end{equation}

With $G_i=\begin{pmatrix}
Id_m & -Id_m & 0 & \hdots &0 \\
Id_m & 0 &-Id_m &\hdots &0 \\
\vdots & & & & \vdots \\
Id_m & 0 & 0 & \hdots & -Id_m
\end{pmatrix}$, $z_i=\begin{pmatrix}
\alpha_i \\
\gamma_{ij} \forall i\sim j
\end{pmatrix}$, $E_i=0$
$\hat H_i= \begin{pmatrix}
H_i & 0 \\
0 & 0 
\end{pmatrix}$ and $\hat F_i^T= \begin{pmatrix}
F_i^T & 0
\end{pmatrix}$, we retrieve and equivalent formulation of (\ref{problème du cours}) in the formalism of the article studied.


The problem statement is to derive an  encrypted ADMM to fulfill the following security requirements:

\begin{enumerate}
    \item Only agent $i$ should have access to $\alpha_i$, $H_i$, $F_i$, $\rho_i$.
    \item Neither the agent nor the operator should have access to plaintext values of $\gamma_i$
    \item Neither the agent nor the operator should have access to plaintext values of $\zeta_i$, tough the agent may know the entries where it is involved.
\end{enumerate}


\subsection{Assumptions}

%Introduce the main assumptions that the paper needs to derive the theoretical results. [1 page max]

Appart the definition of the function for each agent, there is no particular assumptions. We are in the 
good setting of convex, strongly convex (if $H_i$ is positive definite) and smooth function, so linear convergence
holds for peer-to-peer ADMM.


\section{Result: theory and practice}

\subsection{Theory}

% Describe the main theorems of the paper and the results therein in mathematical terms, {\bf by using the notation of the lecture notes.} Interpret the results [1 page max]

Let derive the augmented Lagrangian for each device : 

\begin{align}\label{Lagragian}
L_\rho(z_1,\ldots,z_N,\zeta,\lambda_1,\ldots,\lambda_N)
:=
\sum_{i=1}^{N}
\Big(
\dfrac{1}{2}z_i^TH_iz_i + y_i^TF_iz_i + i_{G_iz_i=0}
+ \lambda_i^{T}\big(z_i - (\zeta)\kappa_i\big)
+ \frac{\rho}{2}\,\| z_i - (\zeta)\kappa_i \|_2^2
\Big)
\end{align}

with the dual variables $\lambda_i \in \mathbb{R}^{md_i}$ and the positive weighting factor $\rho \in \mathbb{R}$. The three ADMM iterations are:

\begin{subequations}
\begin{align}
z_i^{\tau+1}
&:= \arg\min_{z_i}
L_\rho\big(z_1,\ldots,z_M,\zeta^\tau,\lambda_1^\tau,\ldots,\lambda_M^\tau\big), \tag{5a} \\[0.5em]
\zeta^{\tau+1}
&:= \arg\min_{\zeta}
L_\rho\big(z_1^{\tau+1},\ldots,z_M^{\tau+1},\zeta,\lambda_1^\tau,\ldots,\lambda_M^\tau\big), \tag{5b} \\[0.5em]
\lambda_i^{\tau+1}
&:= \lambda_i^\tau + \rho \big( z_i^{\tau+1} - (\zeta^{\tau+1})\kappa_i \big). \tag{5c}
\end{align}
\end{subequations}

The first step (5a) can be distributed, with intermediary quadratic optimization problem under the linear constraint 

\begin{equation*}
    \begin{aligned}
& \min_{z_i} \quad 
\frac{1}{2} z_i^T H_i z_i 
+ y_i^T F_i^T z_i 
+ \lambda_i^{\tau,T} z_i 
+ \frac{\rho}{2} \left\| z_i - (\zeta^\tau)\kappa_i \right\|_2^2 \\
\text{s.t. }& G_i z_i = 0
\end{aligned}
\end{equation*}


and we consequently find $z_i^{\tau+1}$ from solving
\begin{equation}\label{explicit first ADMM step}
\begin{aligned}
\begin{pmatrix}
H_i + \rho I_{v_i} & G_i^T \\
G_i & 0
\end{pmatrix}
\begin{pmatrix}
z_i^{\tau+1} \\
\mu_i^{\tau+1}
\end{pmatrix}
=
\begin{pmatrix}
\rho (\zeta^\tau)\kappa_i - F_i y_i - \lambda_i^\tau \\
0   
\end{pmatrix},
\end{aligned}
\end{equation}

$\mu_i$ values are irrelevant for the algorithm.
After few computation, similar to the point 4.5.1 of the lecture notes, we find for the second step (5b) when $\lambda_i^0=0$ :

\begin{equation}\label{second step ADMM}
    \zeta_i^{\tau+1} = \frac{1}{d_i} \sum_{j \sim i} \gamma_{ji}^{\tau+1}
\end{equation}    
So one can derive the ADMM iterations: \label{ADMM explicit}

\begin{enumerate} \label{ADMM explicit}
    \item  \textbf{Computation step} Each device compute $z_i^{\tau+1}$ thanks to \ref{explicit first ADMM step}
    \item \textbf{Communication step} Each device communicates $\gamma_{ij}^{\tau  +1}$ to its neighbors $j\sim i$. And receive those informations from the neighbors.
    \item \textbf{Computation step}
    Each device comput $\zeta_i$ thanks to \ref{second step ADMM}
    \item \textbf{Communication step}
    Each device communicates $\zeta_{i}^{\tau+1}$ to its neighbors $j\sim i$. And receive those informations from the neighbors.
    \item \textbf{Computation step} Compute $\lambda_i^{\tau+1}$ as in \hyperref[Lagragian]{(5c)}  
\end{enumerate}

\paragraph{Observations}
This algorithm only requires an initial guess for the $(\zeta_i)$, which can be taken as the average of the initial guesses of the neighboring optimal points. The computational steps clearly correspond to ADMM, allowing us to implement it in a distributed manner. 

For the reformulated problem (\ref{main problem}) the previous steps implies that $\forall \tau \quad G_i z_i^\tau = 0 $ so that $z_i$ is just the repetition of $\alpha_i$ several times. That is why in this specific case (the practice case), we can simplify the first computation step with:

\begin{equation*}
\begin{aligned}
& \min_{z_i} \quad 
\frac{1}{2} \alpha_i^T H_i \alpha_i 
+ y_i^T F_i^T \alpha_i 
+ \lambda_i^{\tau,T} (z_i - (\zeta)_{\kappa_i})
+ \frac{\rho}{2} \sum_{j\sim i} \left\| \alpha_i - \zeta_j^\tau \right\|_2^2 \\
\end{aligned}
\end{equation*}
So

\begin{equation}\label{explicit first ADMM step}
(H_i + \rho I_{v_i} )\alpha_i^{\tau+1} =
\rho (\zeta^\tau)\kappa_i - F_i y_i - \lambda_i^\tau 
\end{equation}

the second computation step can be simplified as well:

\begin{equation}\label{second step ADMM}
    \zeta_i^{\tau+1} = \frac{1}{d_i} \sum_{j \sim i} \alpha_{j}^{\tau+1}
\end{equation}   

The article described how homomorphic encryption (HE) is used to implement the ditributed \hyperref[ADMM explicit]{ADMM} in a privacy-preserving way. The original update steps (e.g., $z_i^{\tau+1}$, $\zeta^{\tau+1}$, $\lambda_i^{\tau+1}$) are computed over encrypted data, ensuring that agents cannot access plaintext values. 
Homomorphic encryption allows computations directly on ciphertexts, satisfying
\[
\mathrm{Dec}(\mathrm{Enc}(x)\oplus \mathrm{Enc}(y)) = x + y, \quad
\mathrm{Dec}(\mathrm{Enc}(x)\otimes \mathrm{Enc}(y)) = xy,
\]
which enables secure evaluation of ADMM updates. In practice, a leveled fully homomorphic encryption (FHE) scheme is used to support a limited number of operations without costly bootstrapping. Data is encoded into a finite ring $\mathbb{Z}_q$ via scaling and rounding:
\[
\left\lfloor s \cdot x \right\rceil \bmod q,
\]
ensuring compatibility with encrypted computations. The choice of $q$ balances computational efficiency, numerical precision, and overflow risk.

We denote $[\![ x ]\!] = \mathrm{Enc}_0(x)$.   
Additively homomorphic encryption allows mixed mutliplication between encrypted data $[\![y]\!]_0$ and x:
\begin{equation}
[\![y]\!]_0\odot x = [\![y]\!]_0 \oplus [\![y]\!]_0 ... \oplus [\![y]\!]_0
\end{equation}
many cryptosystem provide an $\odot$ apllication that does not rely on multiple addition. 

We now can adpapt the ADMM setps with encryption
\begin{enumerate} \label{ADMM explicit}
    \item  \textbf{Computation step} Each device compute $[\![z_i^{\tau+1}]\!]_0 $ with:
\begin{equation*}
\begin{aligned}
[\![z_i^{\tau+1}]\!]_0 &= \left( \rho \Gamma_i^{11} \odot [\![(\zeta^\tau)_{\mathcal{K}_i}]\!]_0 \right) \ominus \left( \Gamma_i^{11} \odot [\![\lambda_i^\tau]\!]_0 \right) \\
&\quad \oplus \left( (\Gamma_i^{12} E_i - \Gamma_i^{11} F_i) \odot [\![p_i]\!]_0 \right)
\end{aligned}
\end{equation*}
where 
$
\begin{pmatrix} \Gamma_i^{11} & \Gamma_i^{12} \\ * & * \end{pmatrix} := \begin{pmatrix} H_i + \rho I_{v_i} & G_i^T \\ G_i & 0 \end{pmatrix}^{-1}
$
can be computed once. 
    \item \textbf{Communication step} Each device communicates $[\![\gamma_{ij}^{\tau  +1}]\!]_0$ to its neighbors $j\sim i$. And receive those informations from the neighbors.
    \item \textbf{Computation step}
    Each device comput $[\![\zeta_i]\!]_0$: 
\begin{equation}
    [\![\zeta^{\tau+1}]\!]_0 = \frac{1}{|\mathcal{I}_k|} \odot \left( \bigoplus_{i \in \mathcal{I}_k} [\![(z_i^{\tau+1})_{k_i}]\!]_0 \right)
\end{equation}
    \item \textbf{Communication step}
    Each device communicates $[\![\zeta_{i}^{\tau+1}]\!]_0$ to its neighbors $j\sim i$. And receive those informations from the neighbors.
    \item \textbf{Computation step} Compute $\lambda_i^{\tau+1}$ with:
\begin{equation}
    [\![\lambda_i^{\tau+1}]\!]_0 = [\![\lambda_i^\tau]\!]_0 \oplus \left( \rho \odot \left( [\![z_i^{\tau+1}]\!]_0 \ominus [\![(\zeta^{\tau+1})_{\mathcal{K}_i}]\!]_0 \right) \right)
\end{equation}
\end{enumerate}

At the end of the algorithm, in this scheme, no agent has access to the global secret key $sk_0$. 
Therefore, ciphertexts encrypted under $sk_0$ cannot be directly decrypted by any agent. 
Instead, each agent $i$ owns a personal secret key $sk_i$, ensuring that only agent $i$ can decrypt information intended for it.

To enable the transition from an encryption under $sk_0$ to an encryption under $sk_i$, the protocol relies on a key switching mechanism. An auxiliary key $\mathrm{Enc}_i(sk_0)$ is generated, which corresponds to an encryption of the global secret key $sk_0$ under the public key of agent $i$. 
This information allows the conversion of ciphertexts from the encryption space associated with $sk_0$ to that associated with $sk_i$, without revealing any secret key.

The value $\mathrm{Enc}_i(sk_0)$ is made available to a neighboring agent $j \in \mathcal{N}_i$. 
After $l$ iterations of the algorithm, agent $i$ send $[\![\alpha_i^l]\!]_0$ to agent $j$ which uses this auxiliary key to transform a ciphertext initially encrypted under $sk_0$ into a ciphertext encrypted under $sk_i$. 
As a result, agent $j$ obtains $[\![\alpha_i^l]\!]_i$, which is then transmitted to agent $i$ Since this ciphertext is now encrypted under the key $sk_i$, only agent $i$ can decrypt it and recover the value $\alpha_i^l$.

The encrypted ADMM scheme is designed to meet security objectives under an The primary goals are achieved as follows:
\begin{itemize}
    \item Local parameters $\alpha_i$, $H_i$, $F_i$, and $G_i$ remain confidential because they are only processed in their encrypted forms by neighboring agents.
    \item The privacy of the state variable $\gamma_ij$ is maintained throughout the process, as it is never decrypted during computation.
    \item Although agents process certain entries of the global variable $\zeta$, no individual agent possesses the secret key $sk_0$ required for decryption. 
\end{itemize}
This framework ensures that sensitive local data is protected from unauthorized access during the optimization iterations.

\paragraph{However} the article spot a weakness in process: 

While the scheme is secure under standard assumptions, vulnerabilities may arise if agents deviate from the protocol by sending arbitrary ciphertexts during the final key switch to access information about $sk_0$. To address this, the computation of the local update is outsourced to a neighboring agent. This requires the encryption of the weighting matrices to satisfy security goals, leading to a revised update rule. Consequently, the standard update is replaced by Equation (\ref{new first step}), which leverages homomorphic operations to ensure privacy:

\begin{equation}\label{new first step}
[z_i^{\tau+1}]_0 = ([\rho\Gamma_i^{11}]_0 \otimes [(\zeta^\tau)_{\mathcal{K}_i}]_0) \oplus ([\Gamma_i^{11}]_0 \otimes [\lambda_i^\tau]_0) \oplus ([\Gamma_i^{12} E_i - \Gamma_i^{11} F_i]_0 \otimes [p_i]_0)
\end{equation}

This modification involves encrypted multiplications and slightly increases communication complexity between neighbors, but it fundamentally prevents agents from presenting unintended data during the key switch.

\subsection{Practice}
\label{sec:practice}

%Code the main algorithm of the paper and apply it to the kernel setting of your python project. Compare the result with the most similar algorithm from the course. Comment on the %results. [1 page max] [Send a link/file to the python code].

\paragraph{Kernel Ridge Regression} 
For $5$ agents and $n$ data points, let us define:
\begin{equation}
    f_i(\alpha)=\frac{\sigma^2}{5}\frac{1}{2}\alpha^T K_{mm}\alpha+\frac{1}{2}\lVert y_i-K_{(i)m}\alpha\rVert^2+\frac{\nu}{10}\lVert \alpha\rVert^2.
\end{equation}
Here, $m=\sqrt{n}$ denotes the dimension of $\alpha_i$.

Our problem differs from the one presented in the article, since only agent $i$ contributes to $f_i$. In particular, only $\alpha_i$ appears in the function $f_i$. Therefore, we can define $z_i = \alpha_i$, with no additional information from neighboring agents encoded in $z_i$. Consequently, we introduce a global variable $\zeta \in \mathbb{R}^m$ to enforce consensus.

\textbf{Notation.} In the original problem formulation, $y$ was a vector containing all $n$ target values, and $y_i$ denoted the $i$-th coordinate of $y$ (i.e., a scalar). We modify this notation: here, $y_i$ denotes the subset of data associated with agent $i$ (of size $n/N$, hence a vector). Similarly, $K_{(i)m}$ has size $(n/N)\times m$. 

With this convention, the term 
\[
\frac{1}{2}\sum_{i \in A} (y_i - K_{(i)m}\alpha)^2
\]
(which was incorrectly written in the subject as $\frac{1}{2}\displaystyle\sum_{i \in A}\lVert y_i-K_{(i)m}\rVert^2$) can be properly rewritten as
\[
\frac{1}{2}\lVert y_i-K_{(i)m}\alpha\rVert^2.
\]

The functions $f_i$ can then be expressed as:
\begin{equation}
    f_i(\alpha_i)=\frac{1}{2}\alpha_i^T\left[\frac{\sigma^2}{5}K_{mm}+\frac{\nu}{5}I_m+K_{(i)m}^T K_{(i)m}\right]\alpha_i
    - y_i^T K_{(i)m}\alpha_i - \frac{1}{2}y_i^T y_i.
\end{equation}

The constant term $y_i^T y_i$ can be omitted in the optimization process. 

Using the notation of the article, we define
\[
z_i = \begin{pmatrix} \alpha_i \end{pmatrix},
\]
and impose the global consensus constraint $z_i = \zeta$.

We obtain:
\[
f_i(z_i)=\frac{1}{2}z_i^T H_i z_i - y_i^T K_{(i)m}\alpha_i,
\]
with
\[
H_i=\frac{\sigma^2}{5}K_{mm}+\frac{\nu}{5}I_m+K_{(i)m}^T K_{(i)m}.
\]

\textbf{Important observation.} At this stage, our formulation is very close to a cloud-based version of ADMM (tough the article was focusing on peer-to-peer), where $\zeta$ effectively plays the role of a central server. This is clearly a drawback of our implementation, and a compromise must be made.

The classical ADMM iterations are then:
\begin{align}
z_i^{\tau+1} &= \left( H_i + \rho I \right)^{-1} \left( K_{(i)m}^T y_i + \rho \zeta^\tau - \lambda_i^\tau \right), \\
\zeta^{\tau+1} &= \frac{1}{N} \sum_{i=1}^{N} z_i^{\tau+1}, \\
\lambda_i^{\tau+1} &= \lambda_i^\tau + \rho \left( z_i^{\tau+1} - \zeta^{\tau+1} \right).
\end{align}

The article relies on the \lstinline{OpenFHE} library to implement the CKKS scheme (a fully homomorphic encryption scheme). This library is primarily designed for \lstinline{C++} and is somewhat outdated for Python. Therefore, we instead use the \lstinline{TenSEAL} library as an alternative.

The mapping into a finite ring is treated as a black box: it is convenient to use, although not fully transparent from a theoretical standpoint.

The sketch of our algorithm is the following:
\begin{enumerate}
    \item Initialization of $\lambda_0$ and $\zeta_0$ to $0$. The article states that $\zeta$ should be initialized with a good estimate of $\alpha$; here, we assume that $0$ is a reasonable initial estimate.
    \item Encryption of $\lambda_0$ and $\zeta_0$.
    \item Computation of $z_0$ for all agents in the encrypted domain.
    \item Iteration in the encrypted domain. Note that the update of $\lambda$ is performed in the decrypted domain, using the decrypted value of $\zeta$. This is permissible since each agent $i$ has access to the components it is involved in (this also enable to reduce noise and scaling issues).
\end{enumerate}

We did not manage to implement the key switching, but it was useless here.

The result is illustrated below for the parametr $\rho=1$ and $1000$ iterations (the encrypted algorithm is slower in term of computation time than the classical one).

\begin{figure}[H]
    \centering
    \includegraphics{Figures/GAP_ADMM_FHE.pdf}
    \caption{ADMM with FHE on kernel regression problem}
    \label{fig:my_label}
\end{figure}

\section{Conclusion}
The added value of ADMM-FHE for our problem is quite limited, since each agent has its own objective function that does not depend on the others. As a result, the ADMM scheme obtained using the formulation from the article is very close to a cloud-based implementation, where $\zeta$ effectively plays the role of a central server. This contrasts with the original goal of the article, which was to design an encrypted peer-to-peer ADMM scheme. 

Nevertheless, we implemented encrypted ADMM iterations, and the results are quite promising.

\bibliographystyle{plain} % Définit le style (plain, alpha, abbrv...)
\bibliography{reference}  % Appelle le fichier references.bib

\end{document}



